% 三区/四区稿专用：平台无关的 MIKU 算法主链
\begin{tikzpicture}[
  node distance=8mm and 7mm,
  input/.style={rectangle,rounded corners=2pt,draw=deepblue,fill=inputyellow,
    minimum width=2.6cm,minimum height=.9cm,align=center,font=\footnotesize},
  miku/.style={rectangle,rounded corners=3pt,draw=deeppink,very thick,
    fill=improvedpink!55,minimum width=3.05cm,minimum height=1.05cm,
    align=center,font=\footnotesize},
  solver/.style={rectangle,rounded corners=3pt,draw=deepblue,
    fill=lightblue!50,minimum width=2.65cm,minimum height=1.0cm,
    align=center,font=\footnotesize},
  output/.style={rectangle,rounded corners=2pt,draw=deepblue,fill=outputpink,
    minimum width=2.7cm,minimum height=.9cm,align=center,font=\footnotesize},
  algflow/.style={-Stealth,very thick,deeppink},
  solveflow/.style={-Stealth,thick,deepblue},
  note/.style={font=\scriptsize,align=center,fill=white,inner sep=1pt}
]

\node[input] (input) {道路几何、自车状态\\多障碍物预测};
\node[miku, right=of input] (occupancy) {到达时间查询\\威胁裕度与鲁棒占据};
\node[miku, right=of occupancy] (spatial) {连续分割扫描\\Top-K 空间同伦};
\node[solver, right=of spatial] (pathqp) {路径凸 QP\\求解 $l(s)$};

\node[miku, below=of spatial] (temporal) {安全时间窗\\多冲突时间同伦};
\node[solver, right=of temporal] (speedqp) {双向 ST 边界\\速度凸 QP};
\node[output, right=of speedqp] (trajectory) {时空轨迹\\同伦语义标签};

\draw[algflow] (input) -- (occupancy);
\draw[algflow] (occupancy) -- (spatial);
\draw[solveflow] (spatial) -- node[note,above]{中心路径边界} (pathqp);
\draw[solveflow] (pathqp) |- node[note,pos=.56,right, xshift=2pt]{$l(s)$ 与冲突区} (temporal);
\draw[algflow] (temporal) -- (speedqp);
\draw[solveflow] (speedqp) -- (trajectory);
% 两条反馈边放在图外侧，并将说明文字置于空白区域，避免穿过节点。
\draw[algflow,rounded corners]
  (speedqp.south) -- ++(0,-.55) -- ++(-7.0,0) -- ++(0,1.05)
  node[pos=.18,below, note] {速度解更新 $\tau(s)$} -| (occupancy.south);
\draw[solveflow,rounded corners]
  (trajectory.north) -- ++(0,.45) -- ++(-1.0,0)
  node[pos=.45,above, note] {滚动同伦保持} -| (temporal.east);
\end{tikzpicture}
